% ========== 自动生成的目录内容 ==========
\vspace{8pt}
1. \textbf{寻找凸包（Graham扫描法）}\hfill 2\\
给定平面上的点集Q，使用Graham扫描法求Q的凸包（凸包定义：包含所有点的最小凸多边形，点或在多边形上或在内部），要求展示算法的完整步骤。

\vspace{8pt}
2. \textbf{两个有序数组的中位数}\hfill 2\\
设X[0:n-1]和Y[0:n-1]为两个长度均为n的有序数组，设计O(log n)时间的分治算法，找出X和Y的2n个元素的中位数（偶数个元素取中间两个的平均值，奇数个取最中间元素）。

\vspace{8pt}
3. \textbf{最大子数组问题}\hfill 3\\
给定包含n个整数（有正有负）的数组A，设计O(nlogn)时间的分治算法找出和最大的非空连续子数组，同时设计O(n)时间的Kadane算法。

\vspace{8pt}
4. \textbf{循环移位问题}\hfill 3\\
给定一个按从小到大排序的数组，现将其循环右移若干位（如原数组[1,2,3,4,5]，循环右移2位后为[4,5,1,2,3]），设计算法计算循环右移的位数k（$k \ge 0$），要求时间复杂度尽可能优化。

\vspace{8pt}
5. \textbf{归并排序}\hfill 4\\
给定数组A，使用归并排序算法对其进行升序排序，要求展示分治的"分解$\to$求解$\to$合并"完整过程。

\vspace{8pt}
6. \textbf{折半搜索（查找成功案例）}\hfill 4\\
给定有序数组A，使用折半搜索算法查找目标元素target，要求展示每次查找的low、mid、high值及比较过程，若查找成功返回元素下标，失败返回-1。

\vspace{8pt}
7. \textbf{折半搜索（查找失败案例）}\hfill 5\\
给定有序数组A，使用折半搜索算法查找不存在的目标元素，重点关注查找失败时的边界调整及终止条件判断。

\vspace{8pt}
8. \textbf{大整数乘法（Karatsuba算法）}\hfill 5\\
设X和Y为两个n位大整数（n为偶数），X拆分为A（前n/2位）和B（后n/2位），Y拆分为C（前n/2位）和D（后n/2位），即$X=A \cdot 2^{n/2}+B$，$Y=C \cdot 2^{n/2}+D$。设计优化方案通过减少乘法次数降低时间复杂度。

\vspace{8pt}
9. \textbf{Strassen矩阵乘法}\hfill 6\\
给定两个n$\times$n矩阵A和B，使用Strassen算法（分治法）计算矩阵乘积C=A$\times$B，要求展示7个中间矩阵的定义及复杂度分析。

\columnbreak
\vspace{8pt}
10. \textbf{快速排序}\hfill 6\\
给定数组A，使用快速排序算法对其进行升序排序，要求展示Partition过程及完整的递归排序步骤。

\vspace{8pt}
11. \textbf{第k小元素（选择问题）}\hfill 7\\
给定无序数组A，使用分治法查找该数组的第k小元素，要求展示Partition过程及递归查找步骤。

\vspace{8pt}
12. \textbf{堆排序}\hfill 7\\
给定数组A，使用堆排序算法对其进行升序排序，要求展示最大堆的建立过程、堆顶元素的弹出与堆调整过程。

\vspace{8pt}
13. \textbf{最近点对问题}\hfill 8\\
给定平面上n个点的集合Q，设计分治算法找出距离最近的两个点，要求说明预处理、分解、递归求解、合并的完整流程。

\vspace{8pt}
14. \textbf{天际轮廓问题}\hfill 8\\
给定n座建筑物，每个建筑物用三元组(Li,Ri,Hi)表示（Li为左下x坐标，Ri为右下x坐标，Hi为高度），设计O(nlogn)算法求出这n座建筑物的天际轮廓。

\vspace{8pt}
15. \textbf{两元素和为X的分治算法}\hfill 9\\
给定由n个实数构成的集合S和实数X，判断S中是否存在两个元素的和为X，要求设计分治算法求解。

\vspace{8pt}
16. \textbf{逆序对的分治算法}\hfill 9\\
给定实数序列A，若i<j且A[i]>A[j]，则(i,j)为逆序对。要求用分治方法求序列中逆序对的个数。

\vspace{8pt}
17. \textbf{数组中两元素和大于0的对数量}\hfill 10\\
给定大小为n的数组A，求数组中不同对(i,j)（i<j）的数量，使得A[i]+A[j]>0。

\vspace{8pt}
18. \textbf{分治算法解题的固有方法论}\hfill 11\\
分治算法（Divide and Conquer）的核心思想是"大事化小，小事化了，最后合并"。标准三步走策略：分解(Divide)、解决(Conquer)、合并(Combine)。复杂度分析使用主定理 $T(n) = aT(n/b) + f(n)$。
